% ========================================================
% PROCEDIMIENTO DE LIMPIEZA
% ========================================================
\section{Procedimiento de limpieza del dataset}\label{sec:limpieza}

La limpieza es un pipeline reproducible, ubicado en el directorio \texttt{tp/} del repositorio y
organizado en tres capas (Figura~\ref{fig:pipeline}). \emph{Bronze} guarda los archivos de TLC tal
cual se descargan; \emph{silver}, los viajes que pasan las reglas, con sus columnas originales; y
\emph{gold}, las features en CSV, porque Go no lee Parquet sin librerías de terceros. Polars aplica
las reglas y DuckDB las vuelve a implementar en SQL como auditoría independiente. El procedimiento
también está documentado en \texttt{tp/docs/limpieza.md}.

\begin{figure}[H]
\caption{Pipeline de limpieza por capas}\label{fig:pipeline}
\centering
\small
% Figura: pipeline de limpieza por capas (sección 4).
\begin{tikzpicture}[
    font=\footnotesize,
    node distance=5mm and 20mm,
    caja/.style={draw, rounded corners=1pt, align=center, text width=21mm, minimum height=10mm, inner sep=3pt},
    etiqueta/.style={above, font=\scriptsize},
    flecha/.style={-Stealth, semithick}]
  \node[caja] (tlc) {TLC\\CDN oficial};
  \node[caja, right=of tlc] (bronze) {\textbf{Bronze}\\Parquet original};
  \node[caja, right=of bronze] (silver) {\textbf{Silver}\\Parquet limpio};
  \node[caja, right=of silver] (gold) {\textbf{Gold}\\CSV para Go};
  \node[caja, dashed, below=6mm of silver] (audit) {Auditoría\\SQL en DuckDB};
  \draw[flecha] (tlc) -- node[etiqueta] {descarga} (bronze);
  \draw[flecha] (bronze) -- node[etiqueta] {7 reglas} (silver);
  \draw[flecha] (silver) -- node[etiqueta] {features} (gold);
  \draw[flecha, dashed] (audit) -- (silver);
\end{tikzpicture}

\end{figure}

Todo el proceso se reproduce con tres comandos:

\begin{quote}
\small
\texttt{cd tp \&\& uv sync -{}-extra dev}\\
\texttt{uv run nyc-tlc all \quad\# bronze $\rightarrow$ silver $\rightarrow$ gold y reporte}\\
\texttt{uv run pytest -q \quad\ \ \# pruebas de reglas, features y pipeline}
\end{quote}

\subsection{Reglas de limpieza}

Las siete reglas se aplican en cascada, en el orden de la Tabla~\ref{tab:reglas}, y cada viaje
descartado se atribuye a la primera regla que no cumple. Así, la suma de descartados coincide
exactamente con la diferencia entre bronze y silver. El criterio es quitar lo \emph{imposible}, no
lo \emph{inusual}.

{\small
\begin{xltabular}{\textwidth}{@{}L{1.75cm}L{3.1cm}Yrr@{}}
\caption{Reglas de limpieza aplicadas a enero de 2024}\label{tab:reglas}\\
\toprule
\textbf{Regla} & \textbf{Se queda si} & \textbf{Por qué} & \textbf{Descartados} & \textbf{Auditoría} \\
\midrule
\endfirsthead
\toprule
\textbf{Regla} & \textbf{Se queda si} & \textbf{Por qué} & \textbf{Descartados} & \textbf{Auditoría} \\
\midrule
\endhead
\bottomrule
\endfoot
Requeridos & No hay nulos en inicio, fin, distancia, zonas, tarifa ni total
  & Sin esos campos no se puede validar el viaje ni construir sus features & 0 & 0 \\
\addlinespace
Mes & Empieza en enero de 2024
  & El archivo trae viajes sueltos de otros meses por errores de reloj del taxímetro & 18 & 18 \\
\addlinespace
Duración & Dura entre 60~s y 6~h
  & Menos de un minuto son pruebas o anulaciones; más de seis horas, un taxímetro olvidado
  & 36.893 & 36.893 \\
\addlinespace
Distancia & Mide entre 0 (excluido) y 100 millas
  & Cero indica anulación o GPS sin datos; el máximo crudo, 312.722 millas, es un error de registro
  & 36.284 & 60.430 \\
\addlinespace
Velocidad & Velocidad media $\le$ 80 mph
  & Duración y distancia plausibles por separado pueden ser imposibles juntas & 100 & 1.921 \\
\addlinespace
Montos & Tarifa y total $> 0$
  & Los negativos son reembolsos o anulaciones registrados como filas espejo & 32.272 & 38.341 \\
\addlinespace
Zonas & Origen y destino ubicables
  & Las zonas 264 (\textit{Unknown}) y 265 (\textit{Outside of NYC}) no tienen ubicación
  & 27.571 & 31.527 \\
\midrule
\multicolumn{3}{@{}l}{Bronze (archivo original)} & 2.964.624 & \\
\multicolumn{3}{@{}l}{Total descartado (4,49\,\% de bronze)} & 133.138 & \\
\multicolumn{3}{@{}l}{\textbf{Silver (dataset limpio)}} & \textbf{2.831.486} & 0 \\
\end{xltabular}
{\footnotesize\textit{Nota.} \emph{Descartados} cuenta en cascada, por la primera regla que falla.
\emph{Auditoría} es el conteo independiente de DuckDB sobre bronze, donde un viaje puede violar
varias reglas; por eso algunas cifras son mayores. En silver, las siete reglas dan 0 violaciones.
Las zonas válidas salen de la tabla oficial de TLC, no de una lista escrita a mano.\par}
}

El dataset limpio conserva el 95,5\,\% de los viajes. Tan importante como lo que se descarta es lo
que se conserva a propósito. Los 140.162 viajes con tarifa Flex (\texttt{payment\_type = 0}) traen
nulos en \texttt{RatecodeID}, \texttt{passenger\_count} y \texttt{store\_and\_fwd\_flag}, pero son
viajes reales y ninguna de esas columnas entra al modelo, así que descartarlos borraría el 4,7\,\%
de los datos sin motivo. Por la misma razón se ignoran los errores de captura en campos que no
usamos, como \texttt{passenger\_count = 0} o \texttt{RatecodeID = 99}. No hay duplicados exactos en
el mes, y los extremos plausibles, como un viaje largo a Newark, se amortiguan en gold con escala
logarítmica en lugar de eliminarse.

\subsection{Verificación}

Cada regla tiene pruebas automáticas en sus bordes: el viaje con el valor justo en el umbral se
queda y el siguiente se descarta, sobre registros con el esquema completo de 19 columnas. El
pipeline se prueba además de punta a punta con un bronze en miniatura. La auditoría en DuckDB cierra
el circuito, porque si Polars y DuckDB discreparan, una de las dos implementaciones estaría mal.
Por último, el reporte guarda el sha256 del Parquet de origen (\texttt{c4d59da7bbc8abae\ldots}; el
valor completo está en \texttt{tp/reports/limpieza\_yellow\_2024-01.json}). Como TLC re-publica
meses corregidos, ese hash permite comprobar que otra corrida partió del mismo archivo.

\subsection{Features para K-means}

La capa gold transforma cada viaje en las seis features de la Tabla~\ref{tab:features}. Las zonas
de origen y destino viajan junto a ellas, pero no entran a la distancia: sirven para agregar las
asignaciones por zona después del clustering.

\begin{table}[H]
\caption{Features de la capa gold}\label{tab:features}
\small
\begin{tabularx}{\textwidth}{@{}L{3.5cm}L{4.4cm}Y@{}}
\toprule
\textbf{Feature} & \textbf{Construcción} & \textbf{Por qué} \\
\midrule
\texttt{hora\_sin}, \texttt{hora\_cos} & Hora de inicio en el círculo de 24~h
  & Las 23:00 y la 01:00 quedan a 2 horas, no a 22 \\
\addlinespace
\texttt{dia\_sin}, \texttt{dia\_cos} & Día de la semana en el círculo de 7 días
  & Domingo y lunes quedan como vecinos \\
\addlinespace
\texttt{log\_duracion\_z} & z-score de $\log(1 + \text{minutos})$
  & La cola larga de la duración no domina la distancia euclídea \\
\addlinespace
\texttt{log\_distancia\_z} & z-score de $\log(1 + \text{millas})$ & Mismo motivo que la duración \\
\bottomrule
\end{tabularx}

{\footnotesize\textit{Nota.} Parámetros del escalado: media 2,558 y desviación 0,639 para el
logaritmo de la duración; media 1,170 y desviación 0,656 para el de la distancia. Las features se
revisaron en la PC2 al implementar el lector en Go (Sección~\ref{sec:correcciones}).\par}
\end{table}
